% GENERATED FILE. Edit canonical lessons, then run npm run generate:books.
% canonical-source-hash: fnv1a64:5b006a2af250f023
% canonical-lessons: JA-C439-okubyou, JA-C439-yowamushi, JA-C439-hikyou, JA-C439-ranbou, JA-C439-yuukan

\chapter{Brave and Timid}
\label{ch:ja-brave-and-timid}
\hlchaptermodality{439}

\begin{chapteropening}
\textbf{By the end of this chapter:} \emph{I can describe a brave, timid or rough person.}

Complete the last lesson of chapter 439: I can describe a brave, timid or rough person.
\end{chapteropening}

\section[\texorpdfstring{\ja{おくびょう} (okubyō)}{おくびょう (okubyō)}]{\ja{おくびょう} (okubyō) — timid; cowardly}

\label{lesson:JA-C439-okubyou}



\begin{quote}
Before the new one: say the Japanese for leisurely, then the Japanese for to act big.
\end{quote}

\subsection*{You'll want to know: \ja{おくびょう}}
\textbf{\ja{おくびょう}} — \emph{okubyō} — \textquotedblleft{}timid; cowardly\textquotedblright{}.

\begin{quote}
\textbf{\ja{うちの} \ja{いぬは} \ja{おくびょうです。}} — \emph{uchi no inu wa okubyō desu} — \textquotedblleft{}Our dog is timid.\textquotedblright{}
\end{quote}

A na-adjective: \ja{な} before a noun.

\begin{cousinweb}
\textbf{\ja{おくびょう}} (\emph{okubyō}) is Sino-Japanese (kango): two Chinese characters read with their on'yomi, readings borrowed from Chinese. They are \emph{oku}, \textquotedblleft{}timidity\textquotedblright{} (Middle Chinese \emph{'ik}) and \emph{byō}, \textquotedblleft{}sick\textquotedblright{} (\emph{bjaengH}) (KANJIDIC2; Middle Chinese in Baxter's transcription).
\end{cousinweb}

\begin{grammarlens}[title={What you've built}]
One more word for describing what you lived and felt.
\end{grammarlens}

\subsection*{Your turn}
\begin{itemize}
\raggedright
  \item \emph{Say it:} \emph{okubyō}
  \item \emph{Say it:} \emph{okubyō}, once more
  \item \emph{Say it:} \emph{uchi no inu wa okubyō desu} — Our dog is timid.
  \item \emph{Recall:} say the Japanese for leisurely, then the Japanese for to act big, then say \emph{okubyō} again
\end{itemize}

\begin{tcolorbox}[breakable,colback=teal!4,colframe=teal!35!black,title={Before you move on}]
What does \ja{おくびょう} mean? (\textquotedblleft{}Timid; cowardly\textquotedblright{}.) Say it once more.
\end{tcolorbox}
\section[\texorpdfstring{\ja{よわむし} (yowamushi)}{よわむし (yowamushi)}]{\ja{よわむし} (yowamushi) — a weakling}

\label{lesson:JA-C439-yowamushi}



\begin{quote}
Before the new one: say the Japanese for to act big, then the Japanese for timid.
\end{quote}

\subsection*{You'll want to know: \ja{よわむし}}
\textbf{\ja{よわむし}} — \emph{yowamushi} — \textquotedblleft{}a weakling\textquotedblright{}.

\begin{quote}
\textbf{\ja{よわむしと} \ja{言}\ja{わないで} \ja{ください。}} — \emph{yowamushi to iwanaide kudasai} — \textquotedblleft{}Please do not call me a weakling.\textquotedblright{}
\end{quote}

\begin{cousinweb}
\textbf{\ja{よわむし}} (\emph{yowamushi}) is a native Japanese compound (wago), inherited rather than borrowed from Chinese: \emph{yowa}, \textquotedblleft{}weak (yowai)\textquotedblright{} and \emph{mushi}, \textquotedblleft{}insect\textquotedblright{} (JMdict, KANJIDIC2).
\end{cousinweb}

\begin{grammarlens}[title={What you've built}]
One more word for describing what you lived and felt.
\end{grammarlens}

\subsection*{Your turn}
\begin{itemize}
\raggedright
  \item \emph{Say it:} \emph{yowamushi}
  \item \emph{Say it:} \emph{yowamushi}, once more
  \item \emph{Say it:} \emph{yowamushi to iwanaide kudasai} — Please do not call me a weakling.
  \item \emph{Recall:} say the Japanese for to act big, then the Japanese for timid, then say \emph{yowamushi} again
\end{itemize}

\begin{tcolorbox}[breakable,colback=teal!4,colframe=teal!35!black,title={Before you move on}]
What does \ja{よわむし} mean? (\textquotedblleft{}A weakling\textquotedblright{}.) Say it once more.
\end{tcolorbox}
\section[\texorpdfstring{\ja{ひきょう} (hikyō)}{ひきょう (hikyō)}]{\ja{ひきょう} (hikyō) — unfair; foul}

\label{lesson:JA-C439-hikyou}



\begin{quote}
Before the new one: say the Japanese for timid, then the Japanese for a weakling.
\end{quote}

\subsection*{You'll want to know: \ja{ひきょう}}
\textbf{\ja{ひきょう}} — \emph{hikyō} — \textquotedblleft{}unfair; foul\textquotedblright{}.

\begin{quote}
\textbf{\ja{うそは} \ja{ひきょうです。}} — \emph{uso wa hikyō desu} — \textquotedblleft{}Lying is unfair.\textquotedblright{}
\end{quote}

A na-adjective: \ja{な} before a noun.

\begin{cousinweb}
\textbf{\ja{ひきょう}} (\emph{hikyō}) is Sino-Japanese (kango): two Chinese characters read with their on'yomi, readings borrowed from Chinese. They are \emph{hi}, \textquotedblleft{}lowly\textquotedblright{} (Middle Chinese \emph{pjie}) and \emph{kyō}, \textquotedblleft{}cowardice\textquotedblright{} (\emph{khjaep}) (KANJIDIC2; Middle Chinese in Baxter's transcription).
\end{cousinweb}

\begin{grammarlens}[title={What you've built}]
One more word for describing what you lived and felt.
\end{grammarlens}

\subsection*{Your turn}
\begin{itemize}
\raggedright
  \item \emph{Say it:} \emph{hikyō}
  \item \emph{Say it:} \emph{hikyō}, once more
  \item \emph{Say it:} \emph{uso wa hikyō desu} — Lying is unfair.
  \item \emph{Recall:} say the Japanese for timid, then the Japanese for a weakling, then say \emph{hikyō} again
\end{itemize}

\begin{tcolorbox}[breakable,colback=teal!4,colframe=teal!35!black,title={Before you move on}]
What does \ja{ひきょう} mean? (\textquotedblleft{}Unfair; foul\textquotedblright{}.) Say it once more.
\end{tcolorbox}
\section[\texorpdfstring{\ja{らんぼう} (ranbō)}{らんぼう (ranbō)}]{\ja{らんぼう} (ranbō) — rough; violent}

\label{lesson:JA-C439-ranbou}



\begin{quote}
Before the new one: say the Japanese for a weakling, then the Japanese for unfair.
\end{quote}

\subsection*{You'll want to know: \ja{らんぼう}}
\textbf{\ja{らんぼう}} — \emph{ranbō} — \textquotedblleft{}rough; violent\textquotedblright{}.

\begin{quote}
\textbf{\ja{あには} \ja{うんてんが} \ja{らんぼうです。}} — \emph{ani wa unten ga ranbō desu} — \textquotedblleft{}My elder brother drives roughly.\textquotedblright{}
\end{quote}

A na-adjective: \ja{な} before a noun.

\begin{cousinweb}
\textbf{\ja{らんぼう}} (\emph{ranbō}) is Sino-Japanese (kango): two Chinese characters read with their on'yomi, readings borrowed from Chinese. They are \emph{ran}, \textquotedblleft{}disorder\textquotedblright{} (Middle Chinese \emph{lwanH}) and \emph{bō}, \textquotedblleft{}violence\textquotedblright{} (\emph{bawH}) (KANJIDIC2; Middle Chinese in Baxter's transcription).
\end{cousinweb}

\begin{grammarlens}[title={What you've built}]
One more word for describing what you lived and felt.
\end{grammarlens}

\subsection*{Your turn}
\begin{itemize}
\raggedright
  \item \emph{Say it:} \emph{ranbō}
  \item \emph{Say it:} \emph{ranbō}, once more
  \item \emph{Say it:} \emph{ani wa unten ga ranbō desu} — My elder brother drives roughly.
  \item \emph{Recall:} say the Japanese for a weakling, then the Japanese for unfair, then say \emph{ranbō} again
\end{itemize}

\begin{tcolorbox}[breakable,colback=teal!4,colframe=teal!35!black,title={Before you move on}]
What does \ja{らんぼう} mean? (\textquotedblleft{}Rough; violent\textquotedblright{}.) Say it once more.
\end{tcolorbox}
\section[\texorpdfstring{\ja{ゆうかん} (yūkan)}{ゆうかん (yūkan)}]{\ja{ゆうかん} (yūkan) — courageous; brave}

\label{lesson:JA-C439-yuukan}



\begin{quote}
Before the new one: say the Japanese for unfair, then the Japanese for rough.
\end{quote}

\subsection*{You'll want to know: \ja{ゆうかん}}
\textbf{\ja{ゆうかん}} — \emph{yūkan} — \textquotedblleft{}courageous; brave\textquotedblright{}.

\begin{quote}
\textbf{\ja{おまわりさんは} \ja{ゆうかんでした。}} — \emph{omawarisan wa yūkan deshita} — \textquotedblleft{}The police officer was brave.\textquotedblright{}
\end{quote}

A na-adjective: \ja{な} before a noun.

\begin{cousinweb}
\textbf{\ja{ゆうかん}} (\emph{yūkan}) is Sino-Japanese (kango): two Chinese characters read with their on'yomi, readings borrowed from Chinese. They are \emph{yū}, \textquotedblleft{}courage\textquotedblright{} (Middle Chinese \emph{yowngX}) and \emph{kan}, \textquotedblleft{}daring\textquotedblright{} (\emph{kamX}) (KANJIDIC2; Middle Chinese in Baxter's transcription).
\end{cousinweb}

\begin{grammarlens}[title={What you've built}]
One more word for describing what you lived and felt.
\end{grammarlens}

\subsection*{Your turn}
\begin{itemize}
\raggedright
  \item \emph{Say it:} \emph{yūkan}
  \item \emph{Say it:} \emph{yūkan}, once more
  \item \emph{Say it:} \emph{omawarisan wa yūkan deshita} — The police officer was brave.
  \item \emph{Recall:} say the Japanese for unfair, then the Japanese for rough, then say \emph{yūkan} again
\end{itemize}

\begin{tcolorbox}[breakable,colback=teal!4,colframe=teal!35!black,title={Before you move on}]
What does \ja{ゆうかん} mean? (\textquotedblleft{}Courageous; brave\textquotedblright{}.) Say it once more.
\end{tcolorbox}
